% ============================================================
\chapter{Lebesgue Integral}
\label{ch:integral}
% ============================================================

We arrive at the heart of the edifice: the \emph{Lebesgue integral}. While Riemann sliced the $x$-axis into small intervals and summed the areas of rectangles, Lebesgue had the ingenious idea of slicing the $y$-axis instead: group the points where the function takes similar values, measure these level sets, and sum. This approach, presented in his 1902 thesis, produces an integral that coincides with Riemann's when the latter exists, but applies to an incomparably larger class of functions---and above all, possesses convergence theorems (Fatou, dominated convergence, monotone convergence) of unmatched power.

\section{Integral of simple functions}

\begin{definition}[Integral of a non-negative simple function]
Let $(X, \mathcal{A}, \mu)$ be a measure space and $\varphi = \sum_{k=1}^m c_k
\mathbf{1}_{A_k}$ a non-negative measurable simple function ($c_k \geq 0$). We
define:
\[
\int_X \varphi\, d\mu = \sum_{k=1}^m c_k\, \mu(A_k),
\]
with the convention $0 \cdot (+\infty) = 0$.
\end{definition}

\begin{proposition}[Properties of the integral of simple functions]
\label{prop:simple_integral}
Let $\varphi, \psi$ be non-negative measurable simple functions and
$\alpha, \beta \geq 0$. Then:
\begin{enumerate}
    \item (Linearity) $\int (\alpha\varphi + \beta\psi)\, d\mu =
    \alpha \int \varphi\, d\mu + \beta \int \psi\, d\mu$.
    \item (Monotonicity) If $\varphi \leq \psi$, then
    $\int \varphi\, d\mu \leq \int \psi\, d\mu$.
    \item If $\varphi = 0$ $\mu$-a.e., then $\int \varphi\, d\mu = 0$.
    \item For any $A \in \mathcal{A}$,
    $\int_A \varphi\, d\mu := \int_X \varphi \cdot \mathbf{1}_A\, d\mu$.
\end{enumerate}
\end{proposition}

\begin{proof}
(1) Consider the common partition generated by the level sets of $\varphi$ and
$\psi$. If $\varphi = \sum_i a_i \mathbf{1}_{B_i}$ and $\psi = \sum_j b_j
\mathbf{1}_{C_j}$ with $\{B_i\}$ and $\{C_j\}$ partitions, then the sets
$B_i \cap C_j$ form a finer partition and
\[
\alpha\varphi + \beta\psi = \sum_{i,j} (\alpha a_i + \beta b_j)
\mathbf{1}_{B_i \cap C_j}.
\]
The equality follows by direct computation.

(2) Write $\psi - \varphi \geq 0$ and use linearity.
\end{proof}

\section{Integral of non-negative measurable functions}

\begin{definition}[Integral of a non-negative function]
Let $f : X \to [0, +\infty]$ be measurable. We define:
\[
\int_X f\, d\mu = \sup\left\{\int_X \varphi\, d\mu :
\varphi \text{ simple measurable}, 0 \leq \varphi \leq f\right\}.
\]
\end{definition}

\begin{proposition}[Fundamental properties]
Let $f, g : X \to [0, +\infty]$ be measurable.
\begin{enumerate}
    \item (Monotonicity) If $f \leq g$, then $\int f\, d\mu \leq \int g\, d\mu$.
    \item (Homogeneity) For $c \geq 0$, $\int cf\, d\mu = c \int f\, d\mu$.
    \item $\int f\, d\mu = 0$ if and only if $f = 0$ $\mu$-a.e.
    \item If $\mu(A) = 0$, then $\int_A f\, d\mu = 0$.
\end{enumerate}
\end{proposition}

\begin{proof}
(3) If $f = 0$ a.e., every simple function $0 \leq \varphi \leq f$ satisfies
$\varphi = 0$ a.e., so $\int \varphi\, d\mu = 0$.

Conversely, if $\int f\, d\mu = 0$, set $A_n = \{f > 1/n\}$. Then
$\frac{1}{n} \mathbf{1}_{A_n} \leq f$, so
$\frac{1}{n} \mu(A_n) \leq \int f\, d\mu = 0$, giving $\mu(A_n) = 0$ for all $n$.
Since $\{f > 0\} = \bigcup_n A_n$, we get $\mu(\{f > 0\}) = 0$.
\end{proof}

\begin{theorem}[Chebyshev / Markov inequality]
\label{thm:markov}
Let $f : X \to [0, +\infty]$ be measurable. For every $a > 0$:
\[
\mu(\{f \geq a\}) \leq \frac{1}{a} \int_X f\, d\mu.
\]
\end{theorem}

\begin{proof}
$a \cdot \mathbf{1}_{\{f \geq a\}} \leq f$, so
$a \cdot \mu(\{f \geq a\}) = \int a \cdot \mathbf{1}_{\{f \geq a\}}\, d\mu
\leq \int f\, d\mu$.
\end{proof}

\section{Additivity of the integral (non-negative functions)}

\begin{theorem}
\label{thm:positive_additivity}
Let $f, g : X \to [0, +\infty]$ be measurable. Then:
\[
\int_X (f + g)\, d\mu = \int_X f\, d\mu + \int_X g\, d\mu.
\]
\end{theorem}

This result is a consequence of the monotone convergence theorem
(Chapter~\ref{ch:convergence}). It can also be proved directly by approximation
with simple functions.

\section{Integral of integrable functions}

\begin{definition}[Integrable function]
A measurable function $f : X \to \overline{\R}$ is \emph{integrable} (or
\emph{$\mu$-integrable}, or \emph{summable}) if:
\[
\int_X \abs{f}\, d\mu < \infty.
\]
In this case, we define:
\[
\int_X f\, d\mu = \int_X f^+\, d\mu - \int_X f^-\, d\mu,
\]
where $f^+ = \max(f, 0)$ and $f^- = \max(-f, 0)$. We write
$L^1(X, \mathcal{A}, \mu)$ (or $L^1(\mu)$) for the space of integrable functions.
\end{definition}

\begin{warning}[title=\textbf{The form $\infty - \infty$ is forbidden}]
The integrability condition $\int |f|\, d\mu < \infty$ guarantees that both
$\int f^+\, d\mu$ and $\int f^-\, d\mu$ are finite, thus avoiding the
indeterminate form $\infty - \infty$.
\end{warning}

\begin{theorem}[Properties of the integral]
\label{thm:integral_properties}
Let $f, g \in L^1(\mu)$ and $\alpha, \beta \in \R$.
\begin{enumerate}
    \item (Linearity) $\int (\alpha f + \beta g)\, d\mu =
    \alpha \int f\, d\mu + \beta \int g\, d\mu$.
    \item (Triangle inequality) $\abs{\int f\, d\mu} \leq \int \abs{f}\, d\mu$.
    \item (Monotonicity) If $f \leq g$ a.e., then
    $\int f\, d\mu \leq \int g\, d\mu$.
    \item If $f = g$ a.e., then $\int f\, d\mu = \int g\, d\mu$.
    \item If $f \geq 0$ a.e. and $\int f\, d\mu = 0$, then $f = 0$ a.e.
\end{enumerate}
\end{theorem}

\begin{proof}
(2) We have $-\abs{f} \leq f \leq \abs{f}$, so by monotonicity
$-\int \abs{f}\, d\mu \leq \int f\, d\mu \leq \int \abs{f}\, d\mu$.
\end{proof}

\section{Comparison with the Riemann integral}

\begin{theorem}
Let $f : [a, b] \to \R$ be a bounded function.
\begin{enumerate}
    \item If $f$ is Riemann-integrable, then $f$ is Lebesgue-measurable and the
    two integrals agree:
    \[
    (\mathcal{R})\!\int_a^b f(x)\, dx = \int_{[a,b]} f\, d\lambda.
    \]
    \item $f$ is Riemann-integrable if and only if its set of discontinuities has
    Lebesgue measure zero (Lebesgue's criterion for Riemann integrability).
\end{enumerate}
\end{theorem}

\begin{proof}[Proof idea]
(1) The lower and upper Darboux sums are integrals of simple functions with respect
to $\lambda$. Taking the limit over partitions yields equality.

(2) Let $D$ be the set of discontinuities of $f$. One shows that the upper and
lower Darboux sums converge to the same value if and only if $\lambda(D) = 0$.
\end{proof}

\begin{keyformulas}[title=\textbf{Lebesgue integral --- Summary}]
For $f : X \to [0, +\infty]$ measurable:
\[
\int_X f\, d\mu = \sup\left\{\sum_{k=1}^m c_k \mu(A_k) :
0 \leq \sum_k c_k \mathbf{1}_{A_k} \leq f\right\}.
\]
For general integrable $f$:
\[
\int_X f\, d\mu = \int_X f^+\, d\mu - \int_X f^-\, d\mu.
\]
\end{keyformulas}

\section{The measure $\nu_f$ associated with an integrable function}

\begin{proposition}
If $f \in L^1(\mu)$, $f \geq 0$, the set function
$\nu_f(A) = \int_A f\, d\mu$ is a finite measure on $(X, \mathcal{A})$, absolutely
continuous with respect to $\mu$.
\end{proposition}

\begin{proposition}[Absolute continuity of the integral]
Let $f \in L^1(\mu)$. For every $\varepsilon > 0$, there exists $\delta > 0$ such
that:
\[
\mu(A) < \delta \implies \int_A \abs{f}\, d\mu < \varepsilon.
\]
\end{proposition}

\begin{proof}
Set $f_n = \min(\abs{f}, n)$. Then $f_n \uparrow \abs{f}$ and by monotone
convergence, $\int f_n\, d\mu \uparrow \int \abs{f}\, d\mu < \infty$. Fix $n$ such
that $\int (\abs{f} - f_n)\, d\mu < \varepsilon/2$. Then for $\mu(A) < \delta$
with $\delta = \varepsilon/(2n)$:
\[
\int_A \abs{f}\, d\mu \leq \int_A f_n\, d\mu + \int_A (\abs{f} - f_n)\, d\mu
\leq n \cdot \mu(A) + \varepsilon/2 < \varepsilon. \qedhere
\]
\end{proof}

\section{Exercises}

\begin{exercise}
Compute $\int_{[0,1]} \mathbf{1}_{\Q}\, d\lambda$ and
$\int_{[0,1]} \mathbf{1}_{\R \setminus \Q}\, d\lambda$.
\end{exercise}

\begin{exercise}
Show that if $f \in L^1(\mu)$, then $\{f \neq 0\}$ is $\sigma$-finite (for $\mu$).
\end{exercise}

\begin{exercise}
Let $f : [0, \infty) \to [0, \infty)$ be measurable. Show that
\[
\int_0^\infty f(x)\, dx = \int_0^\infty \lambda(\{f > t\})\, dt.
\]
(\emph{Cavalieri's principle / layer cake formula.})
\end{exercise}

\begin{exercise}
Let $f \in L^1(\R, \lambda)$. Show that
$\lim_{n \to \infty} \int_{\R} f(x) \sin(nx)\, dx = 0$
(\emph{Riemann--Lebesgue lemma}).
\end{exercise}

\begin{exercise}
Let $f \geq 0$ be measurable on $(X, \mathcal{A}, \mu)$. Show that
$\nu(A) = \int_A f\, d\mu$ defines a measure on $(X, \mathcal{A})$ and that
for every measurable $g \geq 0$:
$\int_X g\, d\nu = \int_X gf\, d\mu$.
\end{exercise}

\begin{exercise}
Let $f : \R \to \R$ be Riemann-integrable on every $[a, b]$. Show that $f$ is
Lebesgue-measurable. Is $f$ necessarily in $L^1(\R, \lambda)$?
\end{exercise}
